\documentclass[11pt]{article}
\usepackage[margin=1in]{geometry}
\usepackage{amsmath,amssymb,amsthm}
\usepackage{bm}
\usepackage{xcolor}
\usepackage{hyperref}
\usepackage{enumitem}

\newtheorem{proposition}{Proposition}
\newtheorem{corollary}{Corollary}
\newtheorem{remark}{Remark}
\theoremstyle{definition}
\newtheorem*{note}{Note}

\newcommand{\authorflag}[1]{{\color{purple}\emph{[Draft note: #1]}}}

\title{Round-Indexed Convergence of SCOPE-FD under Partial Participation\\
\large A supplementary derivation for Reviewer Comments R1-C1, R2-C3, and Attached-Review-C1}
\author{Draft prepared for advisor review --- not for submission as-is}
\date{August 2026 (verified against Mu, Garg \& Ratnarajah, \emph{IEEE TCCN}, Vol.~10, No.~4, Aug.~2024)}

\begin{document}
\maketitle

\begin{note}
This note was drafted by an AI assistant. Its first version (an earlier draft
of this file) was built entirely from the manuscript's own citations of
\cite{mu2024fedtskd} and flagged one open question in purple: whether
FedTSKD's Theorem~1 assumes a specific stochastic sampling scheme for the
per-round client subset. \textbf{The source paper has since been supplied and
checked directly, and that question is now resolved} --- see
\S\ref{sec:audit}. Every remaining purple flag below concerns only
manuscript-integration choices (numbering, placement), not the mathematics.
\end{note}

\tableofcontents

% =============================================================================
\section{What FedTSKD actually assumes and proves --- verified against the source}
\label{sec:source-facts}
% =============================================================================

Five facts below are quoted or directly restated from
\cite{mu2024fedtskd}, not inferred from the manuscript's citation of it.

\subsection*{Fact 1: FedTSKD has no client-selection mechanism at all}
Algorithm~1 of \cite{mu2024fedtskd} (``Dynamic Training Steps Based FD
Algorithm'') loops \texttt{for n = 1,\ldots,N do} inside every round $r$ ---
every one of the $N$ clients updates every round. There is no subset, no
selection variable, and no sampling step anywhere in the algorithm. This is
not an oversight to be worked around; it is confirmed as the paper's explicit
experimental design: \emph{``Full participation is considered; that is, all
local users join the training process for each global communication round,
since the communication payload required in the FD algorithm is much less
than that in the FL scheme''} (\cite{mu2024fedtskd}, Section~VI). FedTSKD's
authors chose full participation because FD's communication cost is cheap
enough to afford it at their scale; SCOPE-FD's partial-participation setting
is a genuine extension of the problem they left open, not a departure from an
assumption they made.

\subsection*{Fact 2: Assumptions 1--4 are stated with no reference to a subset}
Quoting \cite{mu2024fedtskd}, Section~II-A, verbatim in substance:
\begin{itemize}[leftmargin=2em]
\item \textbf{Assumption 1 (Smoothness).} $L_1,\ldots,L_N$ are all
$\psi$-smooth.
\item \textbf{Assumption 2 (Convexity).} $L_1,\ldots,L_N$ are all
$\mu$-strongly convex; $Q_1,\ldots,Q_N$ are all $\mu_q$-strongly convex.
\item \textbf{Assumption 3 (Bounded Variance).} The stochastic gradients
satisfy $\mathbb{E}\|\nabla L_n(\bm w_{n,t},\mathcal{D}_{n,t}) - \nabla
L_n(\bm w_{n,t},\mathcal{D}_n)\|_2^2 \le \sigma_n^2$ (and the analogous bound
for $Q_n$ with variance $\sigma_{qn}^2$), where $\mathcal{D}_{n,t}\subset
\mathcal{D}_n$ is \emph{``a subset chosen uniformly and independently from
$n$th user's local training dataset $\mathcal{D}_n$ at time step $t$''} ---
i.e.\ this is a statement about \emph{minibatch sampling within client $n$'s
own data}, not about which clients participate in round $r$.
\item \textbf{Assumption 4 (Uniformly Bounded Gradient).} $\mathbb{E}\|\nabla
L_n(\bm w_{n,t},\mathcal{D}_n)\|_2^2 \le G^2$ (and analogously for $Q_n$ with
bound $G_q^2$).
\end{itemize}
None of the four assumptions mentions $\mathcal{S}_r$, a selection
probability, or any property of \emph{which} clients update in a given round.
They are properties of the loss landscape and of each client's own local
data, full stop.

\subsection*{Fact 3: the proof itself never introduces $\mathcal{S}_r$}
Lemma~1 (local-training step), Lemma~2 (server+local distillation step),
Corollary~1, Corollary~2, and Theorem~1 of \cite{mu2024fedtskd} (its
Appendices B--F) are all stated and proved \emph{per client $n$}, tracking a
single client's own weight sequence $\bm w_{n,t}$ through the recursion
$\bm w_{n,t+1} = \bm w_{n,t} - \eta_t \nabla L_n(\bm w_{n,t},\mathcal{D}_{n,t})$
(their Eq.~12) and its distillation analogue (their Eq.~14). At no point does
any induction step, bound, or constant depend on which \emph{other} clients
did or did not update in the same round. This is exactly what one would
expect given Fact~1: with no selection mechanism in the algorithm, the proof
has nothing selection-related to condition on.

\subsection*{Fact 4: what the time index $t$ actually counts}
$t$ ranges over $\mathcal{T} = \mathcal{T}_{LT} \cup \mathcal{T}_{SD} \cup
\mathcal{T}_{LD}$ (\cite{mu2024fedtskd}, Eq.~11), the merged index set of
local-training, server-distillation, and local-distillation steps. Under
\emph{full} participation this coincides with a single shared wall-clock
counter, since every client passes through the same phases at the same time.
Once participation is partial, this coincidence breaks: a client only
advances through $\mathcal{T}_{LT}$ (and is only handed a server broadcast to
run $\mathcal{T}_{LD}$ on) during rounds in which it is actually selected.
This is the one place a genuinely new statement is needed --- not a violated
assumption, but an index that needs to be correctly re-attached to
per-client participation once selection becomes partial. That is the content
of \S\ref{sec:corollary}.

\subsection*{Fact 5: the objective is already stated per-client, additively}
\cite{mu2024fedtskd}'s own problem statement (their Eq.~2) is
$\bm w^{*} = \arg\min_{\bm w} L_n(\bm w),\ \forall n$ --- $N$ separate,
independent minimizations, one per client, with no shared/global model
coupling them in the stated objective. The manuscript's own
Eq.~(\ref{eq:global_obj}), $F(\{\bm w_n\}) = \sum_n \lambda_n L_n(\bm w_n)$,
inherits this additive structure directly from \cite{mu2024fedtskd} rather
than introducing it. Consequently $F^{*} = \sum_n \lambda_n L_n^{*}$ holds by
the source paper's own problem formulation, not as a simplifying assumption
of this note.

% =============================================================================
\section{The manuscript's given results, restated}
% =============================================================================

\subsection*{Given 1: Proposition 1 (already proved in the manuscript)}
For client $i$ at round $r$, with cumulative participation count
$n_i^{(r-1)}$ and raw debt $d_i^{(r)} = r K/N - n_i^{(r-1)}$, normalized to
$\tilde d_i^{(r)}\in[0,1]$ by min--max scaling, the manuscript's
Proposition~1 proves by induction that whenever $\alpha_u+\alpha_d<1$,
\begin{equation}
\bigl|n_i^{(R)}-n_j^{(R)}\bigr|\;\le\;1 \qquad \forall\, i,j,\ \forall\, R.
\label{eq:given-spread}
\end{equation}

\subsection*{Given 2: the exact participation count (manuscript's Eq.~17)}
Writing $m = KR \bmod N$,
\begin{equation}
n_i^{(R)} \;=\; \Big\lfloor \frac{KR}{N}\Big\rfloor \;+\; \mathbb{1}[i \in \mathcal{M}], \qquad |\mathcal{M}| = m,
\label{eq:given-exact-count}
\end{equation}
so \emph{every} client, not merely the average one, satisfies
\begin{equation}
n_i^{(R)} \;\ge\; \frac{KR}{N} - 1
\qquad \forall\, i.
\label{eq:given-lowerbound}
\end{equation}

\subsection*{Given 3: FedTSKD's per-client bound (verified, \S\ref{sec:source-facts})}
\begin{equation}
\mathbb{E}\bigl[L_n(\bm w_{n,t+1})\bigr] - L_n^{*} \;\le\; \frac{\psi}{2}\cdot\frac{\nu_n}{t+1+\beta_1},
\label{eq:manuscript-conv}
\end{equation}
with $t \in \mathcal{T}_{LT}\cup\mathcal{T}_{SD}\cup\mathcal{T}_{LD}$
counting client $n$'s own elapsed local-training and distillation steps
(Fact~4), and $\psi,\nu_n,\beta_1$ as defined in \cite{mu2024fedtskd}'s
Theorem~1.

% =============================================================================
\section{Why the manuscript's current Remark understates its own case}
% =============================================================================

The manuscript's current Remark (Section~V-B) states that FedTSKD's bound
``continues to hold... without modification.'' Read literally at the level
of Eq.~(\ref{eq:manuscript-conv}), this is close to a tautology, and now that
Facts~1--3 are confirmed it is exactly as trivial as it looks: FedTSKD's
proof was never conditioned on a selection mechanism, because FedTSKD has
none. A generic claim of ``the per-client bound still holds'' therefore
requires no argument at all --- which is precisely Reviewer~1's objection
(Comment~1): the analytical content of the original Remark is minimal because
it was never really about partial participation.

The substantive question the reviewers actually want answered is: as a
function of the number of \emph{communication rounds} $R$, how fast does
every client's bound shrink? Under full participation this is a non-question
(every client's $t$ equals the elapsed round count, up to the fixed
per-round step budget). Under SCOPE-FD's partial participation it is not,
because a client's own clock $t$ only advances on rounds it is selected
(Fact~4), and Eq.~(\ref{eq:manuscript-conv}) alone says nothing about how
small $t$ can be relative to $R$ for an arbitrary policy. This is exactly
where Proposition~1 does real work: it pins down, for \emph{every} client
simultaneously, how $t$ relates to $R$.

% =============================================================================
\section{The round-indexed corollary}
\label{sec:corollary}
% =============================================================================

Let $K_r$ denote the number of local-training steps in round $r$ (as in the
manuscript and in \cite{mu2024fedtskd}; it may be constant or follow a fixed,
data-independent schedule, but it does not depend on the selection policy or
on any client's participation history) and let $S$ denote the fixed number of
distillation steps per round. Define $K_{\min} := \min_{1\le r\le R} K_r > 0$,
a constant fixed by the training-step schedule alone.

\begin{corollary}[Uniform per-client rate in communication rounds]
\label{cor:per-client}
Let Assumptions~1--4 of \cite{mu2024fedtskd} hold (as already assumed by the
manuscript's existing Remark; by Fact~2 none of them constrains the selection
policy), let $\alpha_u+\alpha_d<1$ (as already assumed by Proposition~1), and
let $R$ be large enough that $(K_{\min}+S)\bigl(KR/N - 1\bigr) + \beta_1 > 0$.
Then, under SCOPE-FD selection, \emph{every} client $n\in\{1,\dots,N\}$
satisfies
\begin{equation}
\mathbb{E}\bigl[L_n(\bm w_n^{(R)})\bigr] - L_n^{*}
\;\le\; \frac{\psi}{2}\cdot\frac{\nu_n}{(K_{\min}+S)\!\left(\dfrac{KR}{N} - 1\right) + \beta_1}
\;=\; \mathcal{O}\!\left(\frac{N}{KR}\right),
\label{eq:cor-per-client}
\end{equation}
where $\bm w_n^{(R)}$ denotes client $n$'s local model after $R$
communication rounds.
\end{corollary}

\begin{proof}
By Algorithm~1 of the manuscript (only $n \in \mathcal{S}_r$ update in round
$r$), client $n$'s weights are unchanged on any round it is not selected,
so its own elapsed-step counter $t$ (Fact~4) advances only on rounds it
participates in, and by exactly $K_{r}+S \ge K_{\min}+S$ steps on each such
round. Hence after $R$ communication rounds, client $n$'s own counter
satisfies
\[
t_n^{(R)} \;\ge\; (K_{\min}+S)\cdot n_i^{(R)}.
\]
Applying Eq.~(\ref{eq:manuscript-conv}) at $t+1 = t_n^{(R)}$ and using the
lower bound above together with Eq.~(\ref{eq:given-lowerbound})
(Proposition~1's exact-count refinement, which holds for \emph{every} $i$
simultaneously) gives Eq.~(\ref{eq:cor-per-client}).
\end{proof}

\begin{note}
The word ``every'' is the entire content of this corollary. Proposition~1's
spread guarantee (Eq.~\ref{eq:given-spread}) is what licenses it: no policy
without a comparable balance guarantee --- uniform random selection included
--- can make the same claim uniformly, since an in-expectation fairness
property alone permits some client's realized $n_i^{(R)}$ to fall arbitrarily
below $KR/N$ at any finite $R$.
\end{note}

\begin{corollary}[Global objective, round-indexed]
\label{cor:global}
Under the same hypotheses as Corollary~\ref{cor:per-client}, and using
$F(\{\bm w_n\}) = \sum_n \lambda_n L_n(\bm w_n)$ with $\lambda_n =
|\mathcal{D}_n|/|\mathcal{D}|$, $F^{*}=\sum_n\lambda_n L_n^{*}$ (Fact~5),
\begin{equation}
\mathbb{E}\bigl[F(\{\bm w_n^{(R)}\})\bigr] - F^{*}
\;\le\; \frac{\psi\, N\,\bar\nu}{2\left[(K_{\min}+S)(KR - N) + N\beta_1\right]}
\;=\; \mathcal{O}\!\left(\frac{N}{KR}\right),
\qquad \bar\nu := \sum_{n=1}^N \lambda_n \nu_n.
\label{eq:cor-global}
\end{equation}
\end{corollary}

\begin{proof}
Sum Eq.~(\ref{eq:cor-per-client}) weighted by $\lambda_n$ and use
$\sum_n \lambda_n = 1$ and Fact~5.
\end{proof}

\begin{remark}[Consistency check against full participation]
\label{rem:consistency}
Setting $K = N$ (full participation) in Eq.~(\ref{eq:cor-global}) recovers
$\mathcal{O}(N/(NR)) = \mathcal{O}(1/R)$, matching FedTSKD's original
full-participation rate. This is a necessary consistency check: at $K=N$
every client is selected every round regardless of policy, so SCOPE-FD's
selector has nothing to do, and the bound must degenerate back to
\cite{mu2024fedtskd}'s own rate. It does.

\authorflag{The bash script prepared alongside this note
(\texttt{scripts/run\_convergence\_reference.sh}) runs exactly this
$K{=}N{=}30$ configuration, plus an intermediate $K{=}20$ point, so that
Eq.~(\ref{eq:cor-global})'s predicted $N/(KR)$ scaling can be checked
empirically against the already-available $K \in \{1,3,5,10\}$ points. This
is confirmatory evidence for the reader, not a dependency of the proof above
--- Corollaries~\ref{cor:per-client}--\ref{cor:global} are unconditional
given Assumptions~1--4 and $\alpha_u+\alpha_d<1$, both already assumed
elsewhere in the manuscript.}
\end{remark}

% =============================================================================
\section{Assumption audit --- resolved}
\label{sec:audit}
% =============================================================================

This is the section Reviewer~2 (Comment~3) and the attached review
(Comment~1) explicitly ask for. With \cite{mu2024fedtskd} verified directly
(\S\ref{sec:source-facts}), the audit is no longer conditional:

\begin{enumerate}[label=(\alph*), leftmargin=2em]
\item \textbf{Smoothness (Assumption 1) and convexity (Assumption 2).}
Properties of $L_n$ and $Q_n$ alone (Fact~2). SCOPE-FD changes nothing about
any loss function --- it only changes which clients are handed a training
opportunity in a given round. Transfers unchanged.
\item \textbf{Bounded variance and bounded gradient (Assumptions 3--4).}
Also properties of $L_n$, $Q_n$, and within-client minibatch sampling
(Fact~2), unrelated to the selection of $\mathcal{S}_r$. Transfers unchanged.
\item \textbf{The aggregation rule.} FedTSKD's proof uses $\mathcal{S}_r$
only through the manuscript's own aggregation Eq.~(\ref{eq:agg}); SCOPE-FD
uses the identical rule. Transfers unchanged.
\item \textbf{The sampling-scheme question --- resolved, not merely
audited.} FedTSKD's Theorem~1 is proved under full participation, with no
selection mechanism, no distributional assumption on $\mathcal{S}_r$, and no
step in its proof (Lemmas~1--2, Corollaries~1--2, Appendices B--F of
\cite{mu2024fedtskd}) that references any client other than $n$ itself
(Facts~1 and~3). There is therefore no sampling assumption for a
deterministic rule like SCOPE-FD's to satisfy or violate --- the question the
attached review poses has a clean answer: \emph{all four assumptions
transfer unchanged, because none of them, nor the proof built on them,
was ever about how clients are chosen.} Corollary~\ref{cor:global} is an
unconditional corollary of Theorem~1 of \cite{mu2024fedtskd} composed with
the manuscript's own Proposition~1, not a hedged parallel result.
\end{enumerate}

% =============================================================================
\section{Proposed reply text for the three pending comments}
\label{sec:replies}
% =============================================================================

\subsection*{R1, Comment 1}
\textit{``...the analytical contribution of Section V-B is minimal.''}

\textbf{Reply.} We thank the Reviewer for this comment, and on reflection we
agree that the original Remark understated what needed proving and overstated
what it had shown. The Remark as written was a statement about a single
client's local optimization trajectory, which indeed does not depend on the
selection policy and requires no new argument, exactly as the Reviewer
observes: FedTSKD's Theorem~1 is proved under full participation with no
selection mechanism of any kind, so no step of its proof could have depended
on one. The substantive question is different: how the bound behaves as a
function of communication rounds once participation is partial, since a
client's own local-update clock and the wall-clock round counter are no
longer the same thing once some rounds pass without that client being
selected. We have replaced the Remark with Proposition~\ref{cor:global}
\authorflag{[cite as the manuscript's new Proposition/Corollary number once
inserted]}, which composes FedTSKD's unmodified per-client bound with our own
participation-balance guarantee (Proposition~1) to give a round-indexed rate
that holds uniformly across every client, not merely on average, and which
recovers FedTSKD's original full-participation rate exactly when $K=N$. The
derivation is short, as the Reviewer anticipated --- its value is in making
explicit a composition the original manuscript asserted without stating.

\subsection*{R2, Comment 3}
\textit{``...provide a clearer theorem for the partial-participation setting,
or state more modestly that SCOPE-FD is compatible with the FedTSKD pipeline
under the assumptions of [13].''}

\textbf{Reply.} We thank the Reviewer for offering both paths explicitly. We
have taken the first, and can report that it did not require weakening to the
second: FedTSKD's Assumptions~1--4 are stated entirely in terms of the local
loss functions $L_n,Q_n$ and within-client minibatch sampling, with no
reference to how clients are selected into a round, and its convergence proof
is carried out under full participation with no selection mechanism at all.
Consequently SCOPE-FD's deterministic rotation does not need to be reconciled
with any sampling assumption in [13] --- there isn't one to reconcile with.
Proposition~\ref{cor:global} states, precisely, that composing FedTSKD's
unmodified per-client bound with our own participation-balance guarantee
(already Proposition~1 in the manuscript) gives a round-indexed rate of
$\mathcal{O}(N/(KR))$, which recovers FedTSKD's original $\mathcal{O}(1/R)$
rate exactly at $K=N$ (Remark~\ref{rem:consistency}), verified as a
consistency check rather than assumed.

\subsection*{Attached review, Comment 1}
\textit{``...it would be helpful to clarify whether the assumptions required
in the original convergence analysis of FedTSKD remain satisfied.''}

\textbf{Reply.} We thank the Reviewer for this request and address it
directly. Section~\authorflag{[X]} of the revised manuscript now walks
through FedTSKD's Assumptions~1--4 individually and confirms each transfers
unchanged: the smoothness and convexity assumptions concern only the local
loss functions $L_n$ and $Q_n$; the bounded-variance and bounded-gradient
assumptions concern only within-client minibatch sampling; and the
aggregation-rule step of the proof uses the selected subset only through the
data-size-weighted aggregation, which SCOPE-FD leaves unchanged. Most
directly relevant to the Reviewer's question: FedTSKD's convergence proof
is carried out under full participation and contains no assumption at all
about \emph{how} clients are selected into a round, since the original
algorithm has no such mechanism. There is therefore no sampling assumption
for SCOPE-FD's deterministic rotation to satisfy or violate, and
Section~\authorflag{[X]} states this explicitly rather than leaving it
implicit.

% =============================================================================
\section*{What still needs a human before this goes anywhere}
% =============================================================================
\begin{enumerate}[leftmargin=2em]
\item Decide where Corollary~\ref{cor:global} is placed in Section~V-B of the
manuscript (replacing or following the current Remark), and give it a proper
number; the three reply texts and Remark~\ref{rem:consistency} need that
number substituted in.
\item Confirm whether the manuscript's own $K_r$ (local-training steps per
round) is held constant across the revision's experiments or follows
FedTSKD's original decreasing schedule $K_1\ge K_2\ge\cdots$ --- either way
Corollary~\ref{cor:per-client} holds as stated (it only needs $K_{\min}>0$),
but the manuscript's own text describing $K_r$ should say which.
\item Once \texttt{scripts/run\_convergence\_reference.sh} has been run, add
its five new data points (K=20, K=30=N, five seeds each) to the existing
K-sweep log-log fit and report the empirical exponent alongside
Remark~\ref{rem:consistency}'s predicted $-1$.
\end{enumerate}

\begin{thebibliography}{9}
\bibitem{mu2024fedtskd}
Y.~Mu, N.~Garg, and T.~Ratnarajah, ``Federated Distillation in Massive MIMO
Networks: Dynamic Training, Convergence Analysis, and Communication
Channel-Aware Learning,'' \emph{IEEE Transactions on Cognitive
Communications and Networking}, vol.~10, no.~4, pp.~1535--1550, Aug.~2024.
doi: 10.1109/TCCN.2024.3378215.
\end{thebibliography}

\end{document}
